\documentclass[11pt, a4paper]{article}

\usepackage[T1]{fontenc}
\usepackage[utf8]{inputenc}
\usepackage{tgpagella}
\usepackage[spanish, es-noshorthands]{babel}
\usepackage{amsmath, amssymb, bm}
\usepackage{booktabs}
\usepackage{tabularx}
\usepackage[table, xcdraw]{xcolor}
\usepackage[most]{tcolorbox}
\usepackage[margin=2.5cm]{geometry}
\usepackage{fancyhdr}

\pagestyle{fancy}
\fancyhf{}
\setlength{\headheight}{16pt}
\lhead{\textbf{UCB Sede Tarija} | Ing. Mecatrónica}
\chead{}
\rhead{IMT-121 Circuitos Electrónicos I | Soluciones S06}
\lfoot{Prof: M.Sc. Ing. Bernardo Quiroga Turdera}
\rfoot{Página \thepage}
\renewcommand{\headrulewidth}{0.4pt}

\definecolor{ucbblue}{RGB}{0, 51, 102}

\begin{document}

\begin{tcolorbox}[colback=ucbblue!5!white, colframe=ucbblue, title=\textbf{Solucionario: Teorema de Superposición}]
    \centering \large \textbf{Guía de Resolución y Criterios de Evaluación}
\end{tcolorbox}

\section*{Solución Ejercicio 1}
\begin{itemize}
    \item \textbf{Paso 1 ($V_s$ activa, $I_s$ abierta):}
    El circuito es un divisor de tensión entre $R_1 = 3\,\text{k}\Omega$ y $R_2 = 6\,\text{k}\Omega$.
    \begin{equation*}
        V_o^{(1)} = 9\,\text{V} \cdot \frac{6\,\text{k}}{3\,\text{k} + 6\,\text{k}} = 9 \cdot \frac{6}{9} = \mathbf{6\,\text{V}}
    \end{equation*}
    \item \textbf{Paso 2 ($I_s$ activa, $V_s$ en corto):}
    $V_s = 0\,\text{V}$ se reemplaza por un cortocircuito. Los resistores $R_1$ y $R_2$ quedan en paralelo respecto al nodo $V_o$.
    \begin{equation*}
        R_{\text{eq}} = R_1 \parallel R_2 = \frac{3 \cdot 6}{3 + 6} = 2\,\text{k}\Omega
    \end{equation*}
    Como la fuente $I_s = 1\,\text{mA}$ extrae corriente del nodo $V_o$ hacia tierra, genera una caída de tensión negativa respecto a la referencia positiva del nodo:
    \begin{equation*}
        V_o^{(2)} = -I_s \cdot R_{\text{eq}} = -1\,\text{mA} \cdot 2\,\text{k}\Omega = \mathbf{-2\,\text{V}}
    \end{equation*}
    \item \textbf{Paso 3 (Suma Algebraica):}
    \begin{equation*}
        V_o = V_o^{(1)} + V_o^{(2)} = 6\,\text{V} + (-2\,\text{V}) = \mathbf{4\,\text{V}}
    \end{equation*}
    \item \textbf{Paso 4 (Verificación Nodal):}
    \begin{equation*}
        \frac{V_o - 9}{3\,\text{k}} + \frac{V_o}{6\,\text{k}} + 1\,\text{mA} = 0 \implies 2(V_o - 9) + V_o + 6 = 0 \implies 3V_o = 12 \implies V_o = \mathbf{4\,\text{V}}
    \end{equation*}
    \item \textbf{Paso 5 (Explicación del signo negativo):} La fuente de corriente extrae corriente del nodo de salida, tendiendo a deprimir su potencial frente a tierra, lo cual se opone a la elevación de potencial producida por la fuente de tensión.
\end{itemize}

\section*{Solución Ejercicio 2 (Diagnóstico de Errores)}
\begin{enumerate}
    \item El estudiante sumó magnitudes sin considerar las polaridades relativas ni las direcciones de referencia de las fuentes respecto al nodo de interés.
    \item Si $V_o^{(2)}$ correspondía a una contribución opuesta (como una fuente que extrae corriente o polaridad inversa), se debió restar (sumar algebraicamente con signo negativo). Ignorar esto invalida el resultado.
\end{enumerate}

\section*{Solución Ejercicio 3 (Razonamiento Físico)}
\begin{enumerate}
    \item Una fuente de tensión ideal mantiene una diferencia de potencial fija entre sus terminales independientemente de la corriente. Si su valor es cero, la diferencia de potencial es cero en todo momento, lo cual define exactamente un cortocircuito ideal.
    \item Una fuente de corriente ideal suministra una corriente fija independientemente de la tensión. Si su valor es cero, la corriente impuesta es cero, lo que equivale a una rama interrumpida o circuito abierto.
\end{enumerate}

\section*{Solución Ejercicio 4 (Potencia)}
\begin{enumerate}
    \item $v_{\text{total}} = v_1 + v_2 = 3\,\text{V} + (-1\,\text{V}) = \mathbf{2\,\text{V}}$.
    \item $P_{\text{total}} = \frac{v_{\text{total}}^2}{R} = \frac{2^2}{2000} = \frac{4}{2000} = \mathbf{2\,\text{mW}}$.
    \item $P_1 = \frac{3^2}{2000} = 4.5\,\text{mW}$, $\quad P_2 = \frac{(-1)^2}{2000} = 0.5\,\text{mW}$.
    \quad Suma directa: $P_1 + P_2 = 5\,\text{mW} \neq 2\,\text{mW}$. \\
    \textbf{Explicación:} $P_{\text{total}} = \frac{(v_1+v_2)^2}{R} = \frac{v_1^2 + v_2^2 + 2v_1v_2}{R} = P_1 + P_2 + \frac{2v_1v_2}{R}$. El término cruzado ($2v_1v_2/R = 2(3)(-1)/2000 = -3\,\text{mW}$) explica la diferencia.
\end{enumerate}

\end{document}
